\unit{Exposição 182 — Geometria formal e geometria algébrica}
\sid{182}
\src{193}
\seminar{(Maio de 1959)}
\paperhead{GEOMETRIA FORMAL E GEOMETRIA ALGÉBRICA}%
  {por Alexander GROTHENDIECK}

\fsection{1}{Esquemas}

Sabe-se que um espaço algébrico afim definido sobre um corpo $k$ é
essencialmente determinado por sua álgebra afim $A$ (anel das funções
regulares definidas sobre $k$), pois os morfismos de espaços algébricos
$X\longrightarrow Y$ correspondem biunívocamente aos homomorfismos de
$k$-álgebras $A(Y)\longrightarrow A(X)$. O álgebra afim correspondente a um
espaço algébrico é uma $k$-álgebra de tipo finito e, do ponto de vista
«clássico», não tem elementos nilpotentes; reciprocamente, toda álgebra desse
tipo obtém-se como álgebra afim de um espaço algébrico definido sobre $k$.
Dispõe-se assim de um conhecido dicionário que permite interpretar, em termos
de álgebra comutativa, as situações relativas a espaços algébricos afins.
Há muito se observou que, desse modo, obtêm-se enunciados mais gerais, pois em
geral já não é necessário supor que os anéis em questão sejam do tipo recém
considerado, e a hipótese noetheriana costuma bastar. Em particular, haja ou
não um corpo de base dado, não há razão para excluir o caso em que esses anéis
contenham elementos nilpotentes. Até agora, os geômetras recusaram-se a levar
em conta essas indicações e obstinaram-se em limitar-se a considerar álgebras
afins sem elementos nilpotentes, isto é, espaços algébricos cujos feixes
estruturais não têm elementos nilpotentes (e, na maioria das vezes, até mesmo
espaços algébricos «absolutamente irredutíveis»). O conferencista considera que
essa atitude constituiu um sério obstáculo ao desenvolvimento de métodos
verdadeiramente naturais em geometria algébrica.

Seja $A$ um anel comutativo. É bem sabido que o conjunto
$X=\operatorname{Spec}(A)$ dos ideais primos de $A$ está provido de uma
topologia natural, a «topologia de Zariski» ou topologia espectral. Por outra
parte, há um feixe de anéis comutativos $\Osh_X$ sobre $X$, cujo talo em
$\mathfrak p\in X$ é o anel localizado $A_{\mathfrak p}$ e cujo anel de
seções identifica-se com $A$. Assim, $X$ torna-se um \emph{espaço
anelado}, chamado \emph{espectro primo} de $A$. Um homomorfismo de anéis
$f:A\longrightarrow B$ define um morfismo de espaços anelados
$f':\operatorname{Spec}(B)\longrightarrow\operatorname{Spec}(A)$; a
aplicação subjacente de conjuntos não é outra senão
$\mathfrak p\longmapsto f^{-1}(\mathfrak p)$. Os morfismos

\src{194}
de espaços anelados de $\operatorname{Spec}(B)$ em
$\operatorname{Spec}(A)$ obtidos desse modo são exatamente aqueles para
os quais os homomorfismos $\Osh_x\longrightarrow\Osh_y$ ($x=f'(y)$) são
locais (isto é, a imagem inversa do ideal maximal é o ideal maximal).

Chama-se \emph{esquema afim} um espaço anelado isomorfo a algum
$\operatorname{Spec}(A)$, e \emph{preesquema} um espaço anelado localmente
afim, isto é, tal que cada ponto tem uma vizinhança aberta que, com a
estrutura induzida, é um esquema afim. Os \emph{morfismos} de preesquemas se
definem de maneira evidente; localmente correspondem a homomorfismos de anéis.

Quando se fixa um preesquema $S$ e se consideram morfismos de preesquemas
$X\longrightarrow S$, o espaço $S$ desempenha o papel de um corpo ou de um
anel de base (ou, melhor, de um espaço de base numa fibração). Diz-se
então que $X$ é um $S$-preesquema; se $S=\operatorname{Spec}(A)$, isto
significa também que $\Osh_X$ é um feixe de $A$-álgebras. Assim, todo preesquema
pode ser considerado de maneira única como um $\mathbf Z$-preesquema. Naturalmente,
os $S$-preesquemas formam uma categoria; além disso, demonstra-se que nessa
categoria sempre existe o produto de dois objetos $X,Y$, denotado por
$X\times_S Y$. Essa noção de produto permite definir a mudança de base de um
$S$-preesquema correspondente a um morfismo $S'\longrightarrow S$: com efeito,
$X\times_S S'$ pode ser considerado como um $S'$-preesquema.

Diz-se que $X$ é \emph{separado} sobre $S$ se a diagonal de
$X\times_S X$ é fechada. Chama-se \emph{esquema} um preesquema separado
sobre $\mathbf Z$; ele é então separado sobre qualquer base. Para simplificar,
doravante falaremos apenas de esquemas e, além disso, os suporemos
\emph{noetherianos}, isto é, uniões finitas de abertos afins que são
espectros de anéis noetherianos. Diz-se que $X$ é \emph{de tipo finito
sobre} $S$ se, para todo aberto afim $U$ de $S$, sua imagem inversa em $X$ é
uma união finita de abertos afins cujos anéis são álgebras de tipo finito
sobre o anel de $U$. São esses $S$-esquemas os que se prestam a um estudo
propriamente geométrico. Em particular, para todo $s\in S$, a fibra
$f^{-1}(s)$ de $X$ sobre $s$ é um esquema algébrico sobre o corpo residual
$\kappa(s)$ do anel local $\Osh_s$ de $s$ em $S$. Assim, em certa medida, $X$
pode ser considerado como uma família de «espaços algébricos» $f^{-1}(s)$, onde
o parâmetro $s$ percorre $S$ (isto é, do ponto de vista local, o
conjunto dos ideais primos de um anel dado). Naturalmente, os
$\kappa(s)$ podem ter características distintas. Se
$S=\operatorname{Spec}(k)$, onde $k$ é um corpo, recupera-se essencialmente a
noção usual de «espaço algébrico», com a única diferença de que agora o
feixe estrutural pode ter elementos nilpotentes.

\src{195}
Inspirando-se em noções bem conhecidas, define-se a noção de
\emph{morfismo projetivo} e, mais geralmente, a de \emph{morfismo próprio}.
Tal morfismo é de tipo finito; além disso, transforma partes fechadas em partes
fechadas e conserva essa propriedade sob qualquer mudança de base.

Sendo $X$ um esquema (noetheriano, como sempre), o feixe $\Osh_X$ é um
\emph{feixe coerente de anéis} no sentido de [2]. Portanto, os feixes
coerentes de módulos sobre $X$ são também os feixes que localmente são
isomorfos ao conúcleo de um homomorfismo
$\Osh_X^{m}\longrightarrow\Osh_X^{n}$.

\fsection{2}{Esquemas formais}

Sejam $X$ um esquema e $X'$ uma parte fechada de $X$. Então existe um
subfeixe coerente $J$ de $\Osh_X$ tal que
$X'=\operatorname{Supp}(\Osh_X/J)$ (e inclusive existe um máximo). Provido de
$\Osh_X/J$, o espaço $X'$ torna-se um esquema, denotado por $X_0$;
tal esquema se chama «subesquema fechado de $X$». Para todo $n$ também se
pode considerar $X'$ provido de $\Osh_X/J^{n+1}$, denotado por $X_n$;
é um subpreesquema fechado de $X$ cujo conjunto subjacente continua sendo $X'$,
mas tem outro feixe estrutural, a saber,
$\Osh_{X_n}=\Osh_X/J^{n+1}$. Evidentemente, os $\Osh_{X_n}$ formam um sistema
projetivo de feixes de anéis sobre $X'$, cujo limite projetivo
$\overline{\Osh}_X$ chama-se \emph{completamento formal de $\Osh_X$ ao longo
de $X'$}. Provido desse feixe de anéis, $X'$ chama-se \emph{completamento
formal de $X$ ao longo de $X'$}; é, pois, um espaço anelado, mas em
geral não um esquema. Para todo feixe coerente $F$ sobre $X$, pode-se considerar
analogamente o completamento formal $\overline F=\varprojlim_n F_n$ de $F$ ao
longo de $X'$ (onde $F_n=F\otimes_{\Osh_X}\Osh_X/J^{n+1}$); é um feixe de
módulos sobre $\overline X$. Suas seções chamam-se «seções formais de $F$
ao longo de $X'$» e identificam-se com os elementos de
$\varprojlim_n\Gamma(X',F_n)$. Para $F=\Osh_X$ obtêm-se as «funções
holomorfas» de $X$ ao longo de $X'$ no sentido de Zariski (terminologia que
não seguiremos por suas interferências com a terminologia clássica).

Chama-se \emph{esquema formal} (subentende-se: noetheriano) um espaço
topológico $\mathfrak X$, provido de um feixe de anéis topológicos
$\Osh_{\mathfrak X}$ que satisfaz a condição seguinte: há um isomorfismo
de feixes de anéis topológicos
$\Osh_{\mathfrak X}=\varprojlim_n\Osh_n$, onde os $\Osh_n$ formam um sistema
projetivo de feixes de anéis sobre $\mathfrak X$, cada um dos quais faz
de $\mathfrak X$ um esquema $\mathfrak X_n$, e tal que, para $m\geq n$, o
homomorfismo $\Osh_m\longrightarrow\Osh_n$ é sobrejetivo e tem por núcleo
$J_m^{n+1}$, onde $J_m$ é

\src{196}
o núcleo de $\Osh_m\longrightarrow\Osh_0$. Demonstra-se então que
$\Osh_{\mathfrak X}$ é um \emph{feixe coerente de anéis locais
noetherianos}.

Em virtude das definições, um completamento formal $\overline X$ como o
anterior é um esquema formal e, reciprocamente, todo esquema formal é
\emph{localmente} desse tipo. De fato, dar um esquema formal \emph{afim} (isto
é, tal que $\mathfrak X_0$ seja afim, o que implica que todos os
$\mathfrak X_n$ também o sejam) equivale a dar um anel topológico noetheriano,
separado e completo para a topologia $J$-ádica.

As definições usuais de morfismo, morfismo de tipo finito, morfismo
próprio, etc., para os esquemas ordinários estendem-se sem dificuldade aos
esquemas formais.

\fsection{3}{Os três teoremas fundamentais}

Seja $f:X\longrightarrow Y$ um morfismo próprio de esquemas (noetherianos, como
sempre), seja $Y'$ uma parte fechada de $Y$ e seja $X'$ sua imagem inversa em
$X$; consideremos os completamentos formais correspondentes $\overline Y$ e
$\overline X$. Então $f$ induz um morfismo
$\overline f:\overline X\longrightarrow\overline Y$ de esquemas formais, que
além disso é próprio. Se $F$ é um feixe coerente sobre $X$, então $\overline F$
é um feixe coerente sobre $\overline X$. No teorema 1, esquecem-se $X,Y,F$ e
consideram-se apenas o morfismo próprio $\overline f$ de esquemas formais e o
feixe coerente $\overline F$ sobre $\overline X$. (Não obstante, o conferencista
só escreveu uma demonstração completa para o caso em que se parte de
$X,Y,f,F$.)

\result{TEOREMA}{1 (teorema de finitude)}
\begin{enumerate}
  \item[\textit{i.}] Os $R^{q}\overline f_{*}(\overline F)$ são feixes
    coerentes sobre $\overline Y$.
  \item[\textit{ii.}] Os homomorfismos naturais
    \[
      R^{q}\overline f_{*}(\overline F)
      \longrightarrow \varprojlim_n R^{q}f_{n*}(F_n)
    \]
    são isomorfismos.
\end{enumerate}

Neste enunciado supõe-se escolhido um subfeixe coerente $J$ de $\Osh_Y$ que
define $Y'$; por imagem inversa obtém-se um subfeixe coerente de $\Osh_X$ que
define $X'$ e, por conseguinte, se definem $F_n,X_n,Y_n$ e
$f_n:X_n\longrightarrow Y_n$ como no n\textsuperscript{o} 2. Se partirmos de
um morfismo próprio qualquer entre dois esquemas formais, são evidentes as
pequenas mudanças que devem ser feitas ao explicitar essas notações.

O teorema 1 referia-se apenas à «cohomologia formal». O teorema seguinte

\src{197}
a relaciona com a «cohomologia algébrica» e é análogo a um conhecido teorema
de Serre [4] sobre a comparação entre cohomologia algébrica e cohomologia
analítica.

\result{TEOREMA}{2 (Primeiro teorema de comparação)} Os
$R^{q}f_{*}(F)$ são feixes coerentes sobre $Y$ (o que é um caso particular
do teorema 1), e os homomorfismos naturais
\[
  \overline{R^{q}f_{*}(F)}
  \longrightarrow \varprojlim_n R^{q}f_{n*}(F_n)
\]
são isomorfismos.

\result{COROLÁRIO}{1} Têm-se isomorfismos canônicos:
\[
  \overline{R^{q}f_{*}(F)}=R^{q}\overline f_{*}(\overline F).
\]

Para $q=0$, esse corolário é uma generalização do «teorema fundamental das
funções holomorfas» de Zariski, do qual deduziremos uma generalização do
«teorema de conexão» de Zariski. Observe-se além disso que, enquanto o teorema 1
(\textit{ii}) é trivial para $q=0$, de modo algum ocorre o mesmo com o
teorema 2 nem com sua formulação equivalente (corolário 1). Com efeito, a
demonstração procede por indução descendente sobre $q$ (é trivial para $q$
grande, pois então ambos os membros são nulos), de modo que o caso $q=0$
aparece como a última etapa da indução e, por assim dizer, como o caso «mais
difícil».

\result{COROLÁRIO}{2} Suponhamos $Y=\operatorname{Spec}(A)$ e que $Y'$ está
definido pelo ideal $J$ de $A$. Então, para todo feixe coerente $F$ sobre
$X$, os $H^{q}(X,F)$ são $A$-módulos de tipo finito cujos completamentos
$J$-ádicos são os $H^{q}(\overline X,\overline F)$.

Aplicando, por fim, este corolário a
$H=\Homsh_{\Osh_X}(F,G)$, obtém-se:

\result{COROLÁRIO}{3} Suponhamos $Y=\operatorname{Spec}(A)$ e que $Y'$ está
definido pelo ideal $J$ de $A$. Sejam $F,G$ dois feixes coerentes sobre $X$;
então $\operatorname{Hom}(F,G)$ é um módulo de tipo finito sobre $A$, cujo
completamento $J$-ádico identifica-se com
$\operatorname{Hom}(\overline F,\overline G)$.

Naturalmente, a aplicação natural
$\operatorname{Hom}(F,G)\longrightarrow
\operatorname{Hom}(\overline F,\overline G)$ é a que associa a um
homomorfismo $u:F\longrightarrow G$ sua extensão «por continuidade»
$\overline u:\overline F\longrightarrow\overline G$ (com o que
$\overline F$ se torna um funtor de $F$).

Suponhamos agora que $A$ é separado e completo para sua topologia $J$-ádica;
então os corolários 2 e 3 precedentes dão:
\[
  H^{q}(X,F)=H^{q}(\overline X,\overline F),
  \qquad
  \operatorname{Hom}(F,G)
  =\operatorname{Hom}(\overline F,\overline G).
\]

\src{198}
Esta última identidade mostra que a categoria dos feixes coerentes sobre
$X$ se identifica com uma \emph{subcategoria} (cujos morfismos são os
morfismos induzidos) da categoria dos feixes coerentes sobre
$\overline X$. De fato, tem-se inclusive:

\result{TEOREMA}{3} Para que um feixe de módulos sobre $\overline X$ seja
coerente, é necessário e suficiente que seja isomorfo a um feixe da forma
$\overline F$, onde $F$ é um feixe coerente sobre $X$ (determinado salvo
isomorfismo canônico em virtude do teorema 2, corolário 3). [Lembremos que
agora $Y=\operatorname{Spec}(A)$, sendo $A$ um anel topológico noetheriano,
separado e completo para a topologia $J$-ádica.]

\result{COROLÁRIO}{1} Os subesquemas fechados de $X$ correspondem
biunívocamente aos subesquemas formais fechados de $\overline X$.

Com efeito, correspondem aos subfeixes coerentes de $\Osh_X$ e,
respectivamente, de $\overline{\Osh}_X$. Considerando os grafos de morfismos
como subesquemas fechados, deduz-se do corolário 1:

\result{COROLÁRIO}{2} Sejam $X,Z$ dois esquemas próprios sobre $A$ (anel
noetheriano, separado e completo para a topologia $J$-ádica). Então a
aplicação $g\longmapsto\overline g$ define uma correspondência biunívoca entre
os $Y$-morfismos de $X$ em $Z$ e os $\overline Y$-morfismos de $\overline X$
em $\overline Z$.

Em outros termos, os esquemas algébricos próprios sobre $A$ aparecem como uma
\emph{subcategoria} (cujos morfismos são os morfismos induzidos) da
categoria dos esquemas formais próprios sobre $\overline Y$. No entanto,
deve-se advertir que existem esquemas formais próprios sobre $\overline Y$ que
não são «algebrizáveis», isto é, isomorfos a algum $\overline X$ com $X$
próprio sobre $A$ (do mesmo modo que existem variedades analíticas complexas
compactas que não provêm de variedades algébricas definidas sobre o corpo
dos complexos). Tais esquemas formais aparecem de maneira natural na
«teoria dos módulos». Assinalemos, não obstante, um caso particular interessante em
que um esquema formal é algebrizável:

\result{TEOREMA}{4} Seja $A$ um anel noetheriano local completo de corpo
residual $k$, e seja $\mathfrak X$ um esquema formal próprio sobre $A$ (provido
de sua topologia $\mathfrak m(A)$-ádica). Supõe-se:
\begin{enumerate}
  \item[\textit{i.}] que os anéis locais de $\Osh_{\mathfrak X}$ são
    módulos planos sobre $A$ ou, o que equivale ao mesmo, se se dota a
    $\Osh_{\mathfrak X}$ e a $A$ da filtração definida pelas potências do
    ideal maximal de $A$, para os graduados associados tem-se a relação
    \[
      \operatorname{gr}(\Osh_{\mathfrak X})
      \simeq
      \operatorname{gr}^{0}(\Osh_{\mathfrak X})
      \otimes_k\operatorname{gr}(A).
    \]
  \item[\textit{ii.}] Considerando
    $\mathfrak X_0=\mathfrak X\otimes_A k$ como um esquema algébrico sobre
    $k$, tem-se
    \[
      H^2(\mathfrak X_0,\Osh_{\mathfrak X_0})=0.
    \]
  \item[\textit{iii.}] $\mathfrak X_0$ é projetivo.
\end{enumerate}
Sob essas condições, $\mathfrak X$ é algebrizável e, mais precisamente, é
isomorfo a $\overline X$, onde $X$ é um $A$-esquema projetivo.

As condições \textit{(ii)} e \textit{(iii)} são satisfeitas, em particular,
se $\mathfrak X_0$ é uma curva lisa sobre $k$, e o teorema 4 pode ser aplicado
em particular na «teoria de módulos» para as curvas de gênero dado\ldots{}
Indiquemos como se demonstra o teorema 4: demonstra-se (cf. a proposição 3
mais adiante) que \textit{(i)} e \textit{(ii)} implicam que todo feixe coerente
sobre $\mathfrak X_0$, localmente isomorfo ao feixe fundamental, obtém-se por
redução a partir de um feixe da mesma natureza sobre $\mathfrak X$. Parte-se
então de um feixe «amplo» sobre $\mathfrak X_0$ (existe um por
\textit{(iii)}), levanta-se esse feixe a um feixe invertível sobre $\mathfrak X$ e,
utilizando o teorema 1, se demonstra que um múltiplo deste último define uma
imersão de $\mathfrak X$ no completamento formal de um esquema
$\mathbf P_A^r$ («projetivo tipo» de dimensão $r$ sobre $A$).

Para a demonstração dos teoremas 1 a 3, veja-se [1].

\src{199}
\fsection{4}{Aplicação ao teorema de conexão e ao «Main theorem» de
Zariski}

Seja $f:X\longrightarrow Y$ um morfismo próprio de esquemas. Pelo teorema de
finitude, $f_*(\Osh_X)=\mathcal A$ é um feixe coerente sobre $Y$; além disso, é um
feixe de álgebras comutativas, pelo que corresponde a um $Y$-esquema
$g:Y'\longrightarrow Y$ finito sobre $Y$ (definido pela condição de ser afim
sobre $Y$, isto é, que a imagem inversa de um aberto afim seja afim, e por
$g_*(\Osh_{Y'})=\mathcal A$). É imediato que $f$ se fatora então
canonicamente como $f=gf'$, onde $f':X\longrightarrow Y'$ é um morfismo de
$X$ em $Y'$ que desta vez satisfaz $f'_*(\Osh_X)=\Osh_{Y'}$. Essa fatoração
de $f$ chama-se \emph{fatoração de Stein} de $f$. Aplicando o primeiro
teorema de comparação, corolário 1, a $f'$ e à parte de $Y'$ reduzida a um
ponto $y'$, obtém-se que $f'^{-1}(y')=X'$ é conexo (de outro modo, as
seções formais de $X$ ao longo de $X'$ não formariam um anel local,
enquanto o completamento de
$f'_*(\Osh_X)_{y'}=\Osh_{Y',y'}$ é local!). Demonstramos:

\result{TEOREMA}{5 («Teorema de conexão» de Zariski)} Seja
$f:X\longrightarrow Y$ um morfismo próprio. Então $f$ se fatora de maneira
única (salvo isomorfismo) como $f=gf'$, onde $g:Y'\longrightarrow Y$ é finito
e $f':X\longrightarrow Y'$ satisfaz $f'_*(\Osh_X)=\Osh_{Y'}$ (de onde
$g_*(\Osh_{Y'})=f_*(\Osh_X)$). As fibras de $f'$ são conexas; isto é, o
conjunto das componentes conexas de uma fibra $f^{-1}(y)$ de $f$ está em
correspondência biunívoca com o conjunto dos pontos de $Y'$ sobre $y$, isto
é, com o conjunto dos ideais maximais de $f_*(\Osh_X)_y$.

Daqui se deduzem imediatamente as variantes usuais do teorema de
conexão. Assinalemos aqui apenas o

\src{200}
\result{COROLÁRIO}{1} Para que um ponto $x$ de $X$ esteja isolado em sua fibra
$f^{-1}(y)$, é necessário e suficiente que a fibra $f'^{-1}(y')$ (onde
$y'=f'(x)$) se reduza a $x$, ou que $f'$ induza um isomorfismo de uma
vizinhança de $x$ sobre uma vizinhança de $y'$. O conjunto desses pontos é um
aberto $U$, e $f'$ induz um isomorfismo de $U$ sobre um aberto de $Y'$.

Para demonstrar que $f'$ é um isomorfismo local em $x$, observa-se que $f'$
induz um isomorfismo $\Osh_{y'}\longrightarrow\Osh_x$, como se vê graças a
$f'_*(\Osh_X)=\Osh_{Y'}$ e ao fato de que (sendo $f'$ uma aplicação fechada
cuja fibra em $y'$ se reduz a $x$) os $f'^{-1}(V)$ percorrem um sistema
fundamental de vizinhanças de $x$ quando $V$ recorre um sistema fundamental de
vizinhanças de $y'$. Daqui se deduz imediatamente o resultado seguinte, devido a
Chevalley no caso «geométrico».

\result{COROLÁRIO}{2} Para que $f$ seja um morfismo finito, é necessário e
suficiente que seja próprio e que suas fibras sejam finitas.

Se assim ocorre, $f'$ é efetivamente um isomorfismo pelo que precede.

Seja $f:X\longrightarrow Y$ um morfismo não necessariamente próprio, mas
suponhamos que $X$ está contido como aberto em um $Y$-esquema próprio
$\widetilde f:\widetilde X\longrightarrow Y$ (o que ocorre, em particular,
se $f$ é quase-projetivo). Aplicando o corolário 1, vê-se que $\widetilde f'$
induz um isomorfismo do conjunto $U$ dos pontos de $X$ isolados em sua fibra
sobre um aberto de $Y'$ (e que $U$ é, com efeito, aberto). Deduz-se a
seguinte variante global do «Main Theorem» de Zariski.

\result{TEOREMA}{6 («Main Theorem» de Zariski)} Seja
$f:X\longrightarrow Y$ um morfismo de tipo finito. Então o conjunto $U$ dos
pontos de $X$ que estão isolados em sua fibra é aberto e, se $f$ é
separado,\footnote{A errata de 1962 substitui aqui a hipótese impressa
«$f$ quase-projetivo» pela hipótese mais fraca «$f$ separado», em virtude de
SGA VIII, 6.2: todo morfismo quase-finito e separado é quase-projetivo.} $U$ é
$Y$-isomorfo a uma parte aberta de um esquema $Y'$ finito sobre $Y$.

Como um morfismo de tipo finito é localmente afim e, \textit{a fortiori},
localmente quase-projetivo, do teorema 6 deduzem-se imediatamente as variantes
locais usuais do Main Theorem.

\fsection{5}{Aplicação ao estudo cohomológico dos morfismos próprios e
planos}

Seja $f:X\longrightarrow Y$ um morfismo próprio, e seja $F$ um feixe coerente sobre
$X$; supõe-se que $F$ é $Y$-plano, isto é, os $F_x$ são módulos planos sobre os
anéis $\Osh_y$ ($y=f(x)$). Isto significa também que, para todo $y\in Y$, se
filtrarmos $F$ ao longo da fibra $f^{-1}(y)$ mediante os
$\mathfrak m_y^nF$ (sendo $\mathfrak m_y$ o ideal maximal de $\Osh_y$), o
graduado associado é isomorfo a
$(F/\mathfrak m_yF)\otimes_{\kappa(y)}\operatorname{gr}(\Osh_y)$; em outros
termos, tem-se
\[
  \mathfrak m_y^nF/\mathfrak m_y^{n+1}F
  =F_y\otimes_{\kappa(y)}
  (\mathfrak m_y^n/\mathfrak m_y^{n+1})
\]
para todo inteiro $n$, onde $X_y$ denota a fibra $f^{-1}(y)$ considerada como
esquema próprio sobre o corpo residual $\kappa(y)$ de $y$, e $F_y$ o feixe
$F/\mathfrak m_yF$ induzido por $F$ sobre $X_y$. Tendo em conta esse
isomorfismo e o teorema 2, obtêm-se cotas e, às vezes, cálculos dos
$R^qf_*(F)$ em uma vizinhança de $y$, conhecendo a cohomologia de $X_y$ com
coeficientes em $F_y$. O teorema 2 assume aqui a forma
\[
  \widehat{R^qf_*(F)_y}
  =\varprojlim_n H^q(X_y,F/\mathfrak m_y^nF).
\]
Assinalaremos apenas a consequência seguinte:

\src{201}
\result{PROPOSIÇÃO}{1} Sejam $f:X\longrightarrow Y$ um morfismo próprio e $F$
um feixe sobre $X$, coerente e $Y$-plano. Seja $y\in Y$, seja $q$ um inteiro, e
suponhamos que $H^q(X_y,F_y)=0$. Então $R^qf_*(F)$ é nulo em uma vizinhança de
$y$ e, para todo $n$, o homomorfismo natural
\[
  R^{q-1}f_*(F)_y\longrightarrow
  H^{q-1}(X_y,F/\mathfrak m_y^nF)
\]
é sobrejetivo.

Em particular, se $f$ é um morfismo plano (isto é, $\Osh_X$ é $Y$-plano),
então todo feixe coerente localmente livre $F$ sobre $X$ é $Y$-plano. Sejam
$F,G$ dois feixes desse tipo; aplicando a proposição 1 a
$\Homsh_{\Osh_X}(F,G)$ e a $q=1$, obtém-se:

\result{TEOREMA}{7} Seja $f$ um morfismo próprio e plano, sejam $F,G$ dois feixes
coerentes localmente livres sobre $X$, seja $y\in Y$ e suponhamos que
\[
  H^1\bigl(X_y,\Homsh_{\Osh_{X_y}}(F_y,G_y)\bigr)=0,
\]
então todo homomorfismo $u_0:F_y\longrightarrow G_y$ é induzido por um
homomorfismo
\[
  u:F|V\longrightarrow G|V,
  \qquad V=f^{-1}(U),
\]
onde $U$ é uma vizinhança de $y$.

\result{COROLÁRIO}{} Se $u_0$ é um isomorfismo (resp. um monomorfismo, um
epimorfismo), então o mesmo ocorre com $u$ para $U$ suficientemente
pequeno.

Em particular:

\src{202}
\result{COROLÁRIO}{2} Seja $E_0$ um feixe coerente localmente livre sobre $X_y$
tal que
\[
  H^1\bigl(X_y;\Homsh_{\Osh_{X_y}}(E_0,E_0)\bigr)=0,
\]
então dois feixes localmente livres cujas restrições a $X_y$ sejam isomorfas
a $E_0$ são isomorfos em uma vizinhança de $X_y$.

Assim:

\result{COROLÁRIO}{3} Suponhamos $H^1(X_y,\Osh_{X_y})=0$. Então dois feixes
invertíveis sobre $X$ (isto é, localmente isomorfos a $\Osh_X$) cujas
restrições a $X_y$ sejam isomorfas são isomorfos.

Deduz-se:

\result{PROPOSIÇÃO}{2} Sejam $Y$ um esquema conexo e $E$ um feixe coerente
localmente livre sobre $Y$; consideremos o fibrado de espaços projetivos
$X=\mathbf P(E)$ associado a $E$, provido de seu conhecido feixe invertível
$\Osh_X(1)$. Então todo feixe invertível $L$ sobre $X$ é isomorfo a um feixe da
forma
\[
  f^*(L')\otimes\Osh_X(n),
\]
onde $L'$ é um feixe invertível sobre $Y$ e $n$ um inteiro. Este último está
determinado de maneira única, e $L'$ está determinado salvo isomorfismo.

O corolário 3 precedente demonstra que $L$ é isomorfo a algum $\Osh_X(n)$ em
uma vizinhança de cada fibra. O resto é quase formal.

A proposição 2 permite determinar os $Y$-morfismos de $X=\mathbf P(E)$ em
outro fibrado projetivo. Em particular, obtém-se:

\result{COROLÁRIO}{} Seja $u$ um automorfismo de $X=\mathbf P(E)$. Então
existem um feixe invertível $L'$ sobre $Y$ e um isomorfismo $v$ de $E$ sobre
$E\otimes L'$ tais que $u$ é o isomorfismo correspondente
\[
  \mathbf P(E)\xrightarrow{\sim}\mathbf P(E\otimes L')=\mathbf P(E);
\]
o par $(v,L')$ está determinado salvo isomorfismo.

Seja $\Gamma$ o conjunto das classes de fibrados invertíveis $L'$ sobre $Y$
tais que $E\otimes L'$ seja isomorfo a $E$. Seus elementos são de torção, pois,
se $n$ é o posto de $E$, vê-se (tomando as potências exteriores $n$-ésimas)
que se deve ter $L'^{\otimes n}\simeq\Osh_Y$. O corolário pode ser expresso
então dizendo que se tem uma sequência exata de grupos
\[
  e\longrightarrow
  \operatorname{Aut}(E)/\Gamma(Y,\Osh_Y^*)
  \longrightarrow \operatorname{Aut}_Y(X)
  \longrightarrow \Gamma\longrightarrow e
\]
(deduzida, de resto, da sequência exata de cohomologia obtida a partir
da seguinte sequência exata de feixes de grupos
\[
  e\longrightarrow \Osh_Y^*
  \longrightarrow \underline{\operatorname{Aut}}(E)
  \longrightarrow \underline{\operatorname{Aut}}_Y(X)
  \longrightarrow e,
\]
onde $\Osh_Y^*$ é o feixe das «unidades» de $\Osh_Y$, identificado com o
centro de $\underline{\operatorname{Aut}}(E)$).

\src{203}
\fsection{6}{Aplicação a teoremas de existência e unicidade para feixes e
esquemas sobre um anel completo para a topologia $J$-ádica}

O teorema 7 dava um resultado de unicidade para feixes coerentes localmente
livres utilizando os teoremas 1 e 2. Utilizando o teorema 3, deduzimos agora
teoremas de existência de feixes, de morfismos de esquemas ou de esquemas.
Doravante, $A$ denota um anel noetheriano local, separado e completo. O método
geral consiste sempre em efetuar construções formais, o que consiste
essencialmente em fazer geometria algébrica sobre um anel artiniano e extrair
delas conclusões de natureza «algébrica» mediante os três teoremas
fundamentais.

\result{PROPOSIÇÃO}{3} Seja $\mathfrak X$ um esquema formal próprio e plano
sobre $A$, e seja $F_0$ um feixe localmente livre sobre $X_0$ tal que
\[
  H^2\bigl(X_0,\Homsh_{\Osh_{X_0}}(F_0,F_0)\bigr)=0.
\]
Então existe um feixe localmente livre $F$ sobre $\mathfrak X$ que induz
sobre $X_0$ um feixe isomorfo a $F_0$. (Além disso, este $F$ é único salvo
isomorfismo se
$H^1(X_0,\Homsh_{\Osh_{X_0}}(F_0,F_0))=0$.)

Constroem-se sucessivamente feixes localmente livres $F_n$ sobre os $X_n$, cada
um dos quais induz o anterior. A construção de $F_n$ encontra uma
obstrução em
\[
  H^2\bigl(X_0,\Homsh_{\Osh_{X_0}}(F_0,F_0)\bigr)
  \otimes_{A/J}(J^n/J^{n+1}),
\]
que é, portanto, nula por hipótese. Utilizando agora o teorema 3, obtém-se:

\result{COROLÁRIO}{1} Seja $X$ um esquema próprio e plano sobre $A$, e seja $F_0$
como acima. Então existe um feixe localmente livre $F$ sobre $X$ que induz
sobre $X_0$ um feixe isomorfo a $F_0$. Além disso, este $F$ é único salvo
isomorfismo se
$H^1(X_0,\Homsh_{\Osh_{X_0}}(F_0,F_0))=0$.

Seja $X_0$ um esquema de tipo finito sobre o corpo $k$; supõe-se que $X_0$ é
liso (pelo que entendemos absolutamente liso) sobre $k$, mas não
necessariamente próprio sobre $k$. Seja $A$ um anel artiniano local de corpo
residual $k$. Interessa-nos encontrar os esquemas $X$ planos sobre $A$ tais que
$X\otimes_A k=X_0$ (este é o ponto de partida da «teoria de módulos» ou
das «variações de estrutura» de $X_0$). É equivalente dar um tal $X$, ou dar
sobre o espaço topológico $X_0$ um feixe $\Osh_X$ provido das estruturas
seguintes:
\begin{enumerate}
  \item[\textit{i.}] É um feixe de $A$-álgebras.
  \item[\textit{ii.}] Está provido de um homomorfismo de aumentação
    $\Osh_X\longrightarrow\Osh_{X_0}$ (compatível com as estruturas de
    $A$-álgebras); esses dados estão sujeitos às condições seguintes: a
    aumentação induz um isomorfismo
    $\Osh_X\otimes_A k\xrightarrow{\sim}\Osh_{X_0}$; $\Osh_X$ é plano sobre
    $A$, isto é, o graduado associado a $\Osh_X$, filtrado pelas potências
    do ideal maximal $\mathfrak m$ de $A$, é isomorfo a
    $\operatorname{gr}^0(\Osh_X)\otimes_k\operatorname{gr}(A)$; isto é, se
    têm isomorfismos
    \[
      \mathfrak m^n\Osh_X/\mathfrak m^{n+1}\Osh_X
      =\Osh_{X_0}\otimes_k(\mathfrak m^n/\mathfrak m^{n+1}).
    \]
\end{enumerate}
O fato fundamental é o seguinte:

\src{204}
\result{TEOREMA}{8} Seja $X_0$ um esquema de tipo finito e liso sobre o
corpo $k$, e suponhamos que $X_0$ é afim. Seja $A$ um anel local artiniano de
corpo residual $k$. Então existe um $A$-esquema $X$ plano sobre $A$ tal que
$X\otimes_A k=X_0$, e dois esquemas desse tipo são necessariamente isomorfos.

Observe-se que o isomorfismo em questão não é canônico, pois em geral $X$
terá $A$-automorfismos não triviais que induzem a identidade sobre $X_0$. Por
outro lado, em geral não há uma escolha «canônica» de um $X$ que satisfaça
as condições dadas, salvo quando $A$ é uma $k$-álgebra (caso de
características iguais), em cujo caso pode tomar-se $X=X_0\otimes_k A$, isto
é, $\Osh_X=\Osh_{X_0}\otimes_k A$ (seja ou não afim $X_0$). No caso de
características distintas, ignoro em geral, quando $X_0$ não é afim, se se
pode «levantar» $X_0$ a um $X$ definido sobre $A$. No entanto, seja $n$ um
inteiro $>0$, ponhamos $A_{n-1}=A/\mathfrak m^n$, e suponhamos que se tenha
levantado $X_0$ a um $A_{n-1}$-esquema plano $X_{n-1}$; pretende-se levantar
$X_{n-1}$ a um $A_n$-esquema plano $X_n$. Pelo teorema 8, já se sabe que isso
é possível localmente; por outro lado, verifica-se facilmente que, se $U_n$
levanta um aberto $U_{n-1}$ de $X_{n-1}$, então o feixe dos grupos de
automorfismos de $U_n$ (que induzem a identidade sobre $U_{n-1}$) é
canonicamente isomorfo a
\[
  \Theta_{X_0/k}\otimes_k(\mathfrak m^n/\mathfrak m^{n+1})
\]
restrito a $U_{n-1}$ e, em particular, é comutativo
($\Theta_{X_0/k}$ denota o feixe de germes de $k$-derivações sobre $X_0$).
Segue-se facilmente que há uma obstrução para construir $X_n$ que levante
$X_{n-1}$, a qual se encontra em
\[
  H^2(X_0,\Theta_{X_0/k})\otimes_k
  \mathfrak m^n/\mathfrak m^{n+1}.
\]
Por conseguinte:

\result{COROLÁRIO}{1} Seja $X_0$ um esquema de tipo finito liso sobre $k$, e
suponhamos $H^2(X_0,\Theta_{X_0/k})=0$. Então, para todo anel local
artiniano $A$ de corpo residual $k$, existe um $A$-esquema plano $X$ tal que
$X\otimes_A k=X_0$.

Além disso, se foi possível encontrar um $X$ plano sobre $A$ que levante $X_0$,
então, em virtude do teorema 8, o conjunto das classes (salvo isomorfismo)
de $A$-esquemas planos que levantam $X_0$ identifica-se com
\[
  H^1\bigl(X_0,\underline{\operatorname{Aut}}(X)\bigr),
\]
onde, naturalmente, $\underline{\operatorname{Aut}}(X)$ denota o feixe de
germes de automorfismos do feixe de $A$-álgebras $\Osh_X$ compatíveis com a
aumentação. A filtração de $\Osh_X$ define uma filtração de
$\underline{\operatorname{Aut}}(X)$, e o quociente deste feixe pelo $n$-ésimo
subgrupo da filtração é $\underline{\operatorname{Aut}}(X_n)$; o graduado
associado a essa filtração é comutativo e identifica-se com
$\Theta_{X_0/k}\otimes_k\operatorname{gr}(A)$. Em particular, se
$\mathfrak m^{n+1}$ é a primeira potência de $\mathfrak m$ que é nula,
então $F^n(\underline{\operatorname{Aut}}(X))$ (o último termo da
filtração) pertence ao centro de $\underline{\operatorname{Aut}}(X)$ e é
isomorfo a $\Theta_{X_0/k}\otimes_k\mathfrak m^n$; é também o feixe de
germes de automorfismos de $X$ que induzem a identidade sobre
$X_{n-1}=X\otimes_A A/\mathfrak m^n$. Utilizando estes resultados obtêm-se
imediatamente os enunciados seguintes:

\src{205}
\result{COROLÁRIO}{2} Seja $X_0$ um esquema de tipo finito liso sobre $k$, seja
$A$ um anel local artiniano de corpo residual $k$ e ideal maximal
$\mathfrak m$, suponhamos $\mathfrak m^{n+1}=0$, seja
$A_{n-1}=A/\mathfrak m^n$ e, por fim, seja $X_{n-1}$ um
$A_{n-1}$-esquema plano tal que $X_{n-1}\otimes_A k=X_0$. Então o conjunto
das classes (salvo um isomorfismo que induza a identidade sobre $X_{n-1}$)
de $A$-esquemas planos $X_n$ tais que
$X_n\otimes_A A_{n-1}=X_{n-1}$ é vazio ou é um espaço homogêneo principal
sob
\[
  H^1(X_0,\Theta_{X_0/k})\otimes_k\mathfrak m^n.
\]

(Observe-se que, em geral, não há um elemento origem privilegiado neste
último espaço, pois não há uma maneira privilegiada de levantar $X_{n-1}$ a
$X_n$).

\result{COROLÁRIO}{3} Seja $X_0$ um esquema de tipo finito liso sobre $k$, e
suponhamos $H^1(X_0,\Theta_{X_0/k})=0$. Então, para todo anel local
artiniano $A$ de corpo residual $k$, existe no máximo (salvo isomorfismo) um
$A$-esquema plano $X$ tal que $X\otimes_A k=X_0$.

Os corolários 1 e 3 implicam imediatamente os enunciados, de aparência mais
geral, que se obtêm supondo apenas que $A$ é um anel local
noetheriano completo de corpo residual $k$, contanto que se introduza aí $X$
como um esquema formal sobre $A$:

\result{TEOREMA}{9} Sejam $k$ um corpo e $X_0$ um esquema de tipo finito liso
sobre $k$. Para todo anel local noetheriano completo $A$ de corpo residual
$k$, seja $F(A)$ o conjunto das classes (salvo um isomorfismo que induza a
identidade sobre $X_0$) de esquemas formais $\mathfrak X$ sobre $A$, de tipo
finito e planos sobre $A$, tais que $\mathfrak X\otimes_A k=X_0$. Com estas
notações:
\begin{enumerate}
  \item[\textit{i.}] Se $H^1(X_0,\Theta_{X_0/k})=0$, então, para todo $A$,
    $F(A)$ tem no máximo um elemento.
  \item[\textit{ii.}] Se $H^2(X_0,\Theta_{X_0/k})=0$, então, para todo $A$,
    $F(A)$ tem ao menos um elemento.
\end{enumerate}

\src{206}
\result{COROLÁRIO}{1} Suponhamos que $X_0$ é próprio sobre $k$. Sob a
condição \textit{(i)}, para todo $A$ existe no máximo (salvo um isomorfismo que
induza a identidade sobre $X_0$) um esquema $X$ próprio e plano sobre $A$ tal
que $X\otimes_A k=X_0$.

Utiliza-se o teorema 3, corolário 2. Por exemplo:

\result{COROLÁRIO}{2} Se $X$ é um esquema próprio e plano sobre $A$ tal que
$X\otimes_A k$ é isomorfo ao esquema projetivo tipo $\mathbf P_k^r$ de
dimensão $r$ sobre $k$, então $X$ é isomorfo a $\mathbf P_A^r$.

(Também se pode deduzir esse resultado da proposição 3, corolário 1).

\result{COROLÁRIO}{3} Seja $X_0$ um esquema liso e projetivo sobre $k$, e
suponhamos que se tenha
\[
  H^2(X_0,\Osh_{X_0})=H^2(X_0,\Theta_{X_0/k})=0.
\]
Então, para todo $A$, existe um esquema projetivo e plano $X$ sobre $A$ tal
que $X\otimes_A k=X_0$.

Combinam-se o teorema 9\textit{(ii)} e o teorema 4. Em particular:

\result{COROLÁRIO}{4} Seja $X_0$ o esquema de uma curva algébrica completa e
lisa sobre $k$. Então, para todo anel local noetheriano completo $A$ de
corpo residual $k$, existe um «esquema de curva lisa» $X$ sobre $A$ tal que
$X\otimes_A k=X_0$.

\result{OBSERVAÇÕES}{}

1) Os corolários 3 e 4 são especialmente interessantes se $k$ é de
característica $p\ne0$, tomando como $A$ um anel de valoração discreta de
característica $0$ e corpo residual $k$; por exemplo, o «menor $A$ possível»,
a saber, aquele cujo ideal maximal é gerado por $p$. (De fato, pelos
teoremas de Cohen, basta conhecer os corolários 3 e 4 para um anel $A$ desse
tipo). Assinalemos a este respeito que, segundo os especialistas, não se sabe se
existem esquemas sobre um corpo $k$ que não sejam a redução módulo $p$ de um
esquema plano definido sobre um tal anel $A$. Ao menos, os resultados deste
número proporcionam um meio de investigar sistematicamente essa questão. Dever-se-ia
começar por examinar se a primeira obstrução que aparece em
$H^2(X_0,\Theta_{X_0/k})$ é necessariamente nula.

\emph{Complemento de 1962.} --- J.-P. Serre construiu uma variedade
projetiva não singular, de dimensão 3, sobre um corpo algebricamente fechado
$k$ de característica $p>0$, que não se obtém por redução de um esquema
próprio sobre um anel local íntegro de corpo residual $k$ e corpo de
frações de característica zero. Mumford teria encontrado um exemplo análogo
com uma superfície projetiva não singular.

2) Observe-se que o teorema 3, e a técnica correspondente, só são válidos para
um anel de base completo (local, para fixar as ideias).

Para passar de resultados conhecidos para o completamento de um anel local aos
resultados correspondentes para esse mesmo anel local, seria necessário um
quarto «teorema fundamental», cujo enunciado definitivo ainda está por encontrar.

\src{207}
3) Comparem-se os resultados deste número (em particular, os corolários 1 e
2 precedentes) e do seguinte com os resultados de Kodaira--Spencer sobre a
variação das estruturas complexas. Utilizando o teorema conjectural ao qual
se acaba de aludir, sob as condições do corolário 1, mas sem supor já
que $A$ seja completo, se deveria poder concluir que existe um anel $A'$ que
contém $A$, finito e não ramificado sobre $A$, tal que $X\otimes_A A'$ e
$X'\otimes_A A'$ são $A'$-isomorfos (sendo $X,X'$ dois $A$-esquemas próprios e
planos dados, tais que $X\otimes_A k=X'\otimes_A k=X_0$). Isto é ao menos o
que pode demonstrar-se quando $X_0=\mathbf P_k^r$, utilizando o corolário da
proposição 2. Em todo caso, quando $H^1(X_0,\Theta_{X_0/k})=0$, pode
demonstrar-se que as fibras de $X,X'$ sobre todo ponto $y$ de
$Y=\operatorname{Spec}(A)$ são isomorfas, ao menos depois de passar ao
fecho algébrico do corpo residual $\kappa(y)$. (Tem-se um resultado
local, aparentemente mais forte, sem supor já necessariamente que $A$ seja
local). Quanto às «variações de estrutura» do espaço projetivo,
assinalemos além disso a questão seguinte, sugerida por um problema correspondente
de Kodaira--Spencer. Seja $X$ um esquema próprio e plano sobre o anel local
íntegro $A$ de corpo de frações $K$ e corpo residual $k$, e suponhamos que
$X\otimes_A K$ é isomorfo a $\mathbf P_K^r$; é verdade que
$X\otimes_A k=X_0$ é isomorfo a $\mathbf P_k^r$ (ao menos sobre o fecho
algébrico de $k$)? Nesse problema pode-se supor que $A$ é um anel de
valoração discreta completo. Coloca-se um problema análogo quando $X_0$ é
uma variedade abeliana.

\emph{Complemento de 1962.} --- O conferencista é agora menos otimista
com respeito aos resultados conjecturados acima. Não obstante, Hironaka resolveu
afirmativamente o problema relativo às variações de estrutura
do espaço projetivo, e Koizumi o problema análogo para as variedades
abelianas.

\fsection{7}{Aplicação à «Teoria dos módulos»}

Como o próprio conferencista apenas começou a abordar essa teoria, seremos
obrigados a nos limitar a algumas indicações. Para simplificar, situamo-nos
sobre um corpo $k$, isto é, trabalhamos em igual característica, embora o
teorema 8 permita tratar também o caso geral sem mudança essencial, ao que
parece. Atualmente não superamos a etapa «formal»; no entanto, o
conferencista espera poder construir a partir dela verdadeiros esquemas de
módulos em certos casos e, em particular, construir para cada inteiro $g$ um
esquema sobre os inteiros que desempenhe o papel de esquema de módulos universal
para as curvas lisas de gênero $g$.

\emph{Complemento de 1962.} --- Mumford construiu os esquemas de módulos
para as curvas de gênero $g$ (cf. SHMI). O teorema 10 demonstra além disso que
os «esquemas de Jacobi de nível $n$» da teoria de módulos são não singulares,
e inclusive lisos sobre $\mathbf Z$.

\src{208}
Retomemos a situação e as notações do teorema 9, supondo agora que $A$
é uma álgebra local de posto finito sobre $k$, que para simplificar se supõe
algebricamente fechado. Então $F(A)$ pode ser considerado como um funtor
covariante de $A$, com valores na categoria de conjuntos: um homomorfismo de
$k$-álgebras $A\longrightarrow B$ define uma aplicação
$F(A)\longrightarrow F(B)$, pois todo $A$-esquema plano $X$ tal que
$X\otimes_A k=X_0$ dá lugar a um $B$-esquema $X\otimes_A B$ com as mesmas
propriedades. Suponhamos que se possa encontrar uma $k$-álgebra local noetheriana
completa $\mathcal O$ e um isomorfismo funtorial
\begin{equation}
  \tag{*}
  \operatorname{Hom}(\mathcal O,A)\xrightarrow{\sim}F(A)
\end{equation}
(onde o primeiro membro denota os homomorfismos de $k$-álgebras). Vê-se
facilmente que um tal $\mathcal O$ está determinado salvo isomorfismo canônico;
chama-se então ao espectro formal $\mathfrak Y$ de $\mathcal O$ (isto é,
ao espaço topológico reduzido a um ponto e provido de um feixe de anéis
topológicos reduzido a $\mathcal O$) o \emph{esquema formal de módulos} para
$X_0$. (Observe-se que ele não existe necessariamente). Seja $\mathfrak r$ o ideal
maximal de $\mathcal O$; para cada $n$, seja
$\mathcal O_n=\mathcal O/\mathfrak r^{n+1}$ (portanto $\mathcal O_0=k$).
Então o homomorfismo canônico $\mathcal O\longrightarrow\mathcal O_n$ é um
elemento de $\operatorname{Hom}(\mathcal O,\mathcal O_n)$, pelo que define um
elemento de $F(\mathcal O_n)$, isto é, um $\mathcal O_n$-esquema plano $X_n$
cuja redução módulo $\mathfrak r$ é $X_0$. Esses $X_n$ se deduzem uns dos
outros por extensão de escalares (isto é, aqui por reduções), de onde
resulta que provêm de um esquema formal $\mathfrak X$ bem determinado sobre
o esquema formal de módulos $\mathfrak Y$; $\mathfrak X$ é plano sobre
$\mathfrak Y$ e $\mathfrak X_0=X_0$. Como se vê imediatamente, o isomorfismo (*)
obtém-se então associando a todo homomorfismo de $k$-álgebras
$\mathcal O\longrightarrow A$ a classe do $A$-esquema
$\mathfrak X\otimes_{\mathcal O}A$ (isto é, a todo morfismo de $k$-esquemas
$Y'=\operatorname{Spec}(A)\longrightarrow\mathfrak Y$ associa-se o
$Y'$-esquema $\mathfrak X\otimes_{\mathfrak Y}Y'$ obtido por mudança de base).
Além disso, vê-se que o isomorfismo (*) e sua descrição precedente continuam
válidos se somente se supõe que $A$ é uma $k$-álgebra local noetheriana e
completa (não necessariamente artiniana). Naturalmente, como de costume,
$\mathcal O$ pode perfeitamente ter \textit{a priori} elementos nilpotentes,
e parece provável que existam casos em que $\mathcal O$ seja ele mesmo artiniano
sem ser idêntico a $k$. Isto mostra até que ponto o ponto de vista de
Kodaira--Spencer (que se limita a tomar como $A$ anéis regulares) é
inadequado \textit{a priori} no caso geral.

\src{209}
Fica por dar condições suficientes para que exista um esquema formal de
módulos para $X_0$, que se supõe próprio sobre $k$. Em geral, é fácil dar
condições necessárias e suficientes simples sobre um funtor
$A\longmapsto F(A)$ (de $k$-álgebras locais de posto finito, com valores nos
conjuntos) para que possa ser escrito na forma
$\operatorname{Hom}(\mathcal O,A)$ para um $\mathcal O$ adequado. Não as
detalharemos aqui. Assinalemos apenas que, no caso que nos ocupa, essas
condições impõem sobre $X_0$ condições cohomológicas não triviais, e parece
pouco provável que sejam sempre satisfeitas, embora o conferencista não tenha
construído nenhum contraexemplo. Em contrapartida, parece plausível que a condição
$H^0(X_0,\Theta_{X_0/k})=0$ seja uma condição suficiente (embora de modo algum
necessária) para a existência de um esquema formal de módulos. Limitamo-nos a
enunciar aqui um teorema num caso particularmente simples (cujo análogo na
teoria de espaços analíticos é bem conhecido, cf. Kodaira--Spencer), que se
estabelece sem dificuldade mediante os resultados do número anterior:

\result{TEOREMA}{10} Seja $X_0$ um esquema próprio e liso sobre o corpo $k$
tal que
\[
  H^0(X_0,\Theta_{X_0/k})=H^2(X_0,\Theta_{X_0/k})=0.
\]
Então existe um esquema formal de módulos para $X_0$, correspondente a um
anel local regular $\mathcal O$ (isto é, uma álgebra de séries formais sobre
$k$).

Como já assinalamos, não é certo em geral que o esquema formal
$\mathfrak X$ sobre $\mathcal O$ seja algebrizável; mas sabe-se que o é
quando $X_0$ é projetivo e $H^2(X_0,\Osh_{X_0})=0$ (teorema 4), por exemplo
quando $X_0$ tem dimensão 1. Isto dá certo fundamento à esperança de
construir um esquema de módulos sobre os inteiros para as curvas de gênero
dado\ldots{}

Observe-se também que métodos como os expostos neste número podem ser aplicados
à construção e ao estudo das variedades de Picard, assim como a muitas
outras construções. Voltaremos a isso em breve.

\fsection{8}{Aplicação ao grupo fundamental}

As técnicas expostas permitem abordar o estudo sistemático do grupo
fundamental seguindo o modelo da teoria topológica. Os dois primeiros
teoremas enunciados neste número são generalizações de resultados de um
trabalho recente de Lang--Serre.

Seja $X$ um esquema. Um $X$-esquema $X'$ se chama recobrimento não ramificado de
$X$ se
\begin{enumerate}
  \item[\textit{i.}] $X'$ é finito sobre $X$, isto é, está definido por um
    feixe coerente de álgebras $\mathcal A=\mathcal A(X')$ sobre $X$;
  \item[\textit{ii.}] $\mathcal A$ é um feixe localmente livre sobre $X$;
  \item[\textit{iii.}] Para todo $x\in X$, o quociente
    \[
      \mathcal A_x/\mathfrak m_x\mathcal A_x
      =\mathcal A_x\otimes_{\Osh_x}\kappa(x)
    \]
    é um álgebra separável sobre $\kappa(x)$.
\end{enumerate}

\src{210}
Essa noção de recobrimento não ramificado (devida a Serre e ao conferencista)
possui todas as propriedades elementares que se pode razoavelmente esperar, cuja
lista não daremos. Limitemo-nos a dizer que dá lugar a uma teoria de Galois
modelada sobre a teoria de Galois clássica (e que a contém; as
demonstrações são bem mais simples que as comumente aceitas para
esta última) e sobre a teoria de Galois dos recobrimentos topológicos. Mais
precisamente, chamemos \emph{ponto geométrico} de um esquema $X$ um morfismo
$a$ do espectro $\xi$ de um corpo algebricamente fechado $\Omega$ em $X$,
isto é, o dado de uma extensão algebricamente fechada do corpo residual
$\kappa(x)$ de um ponto $x=|a|$ de $X$ (chamado \emph{localização} do ponto
geométrico $a$). Se $X'$ é então um recobrimento não ramificado de $X$, se
lhe pode associar o conjunto $E_a(X')$ dos «pontos geométricos de $X'$ sobre
$a$», isto é, o conjunto dos pares formados por um $x'\in X'$ sobre $x$
e um $\kappa(x)$-homomorfismo $\kappa(x')\longrightarrow\Omega$. Obtém-se assim
(para $(X,a)$ fixos) um funtor $F(X,a)$ da categoria $R(X)$ dos
recobrimentos não ramificados $X'$ de $X$ na categoria dos conjuntos
finitos. Se $X$ é conexo, o par formado por $R(X)$ e $F(X,a)$ possui as
propriedades formais necessárias para ser isomorfo ao par análogo definido por
um grupo topológico compacto totalmente descontínuo $\pi$ adequado (isto é,
um limite projetivo de grupos finitos): toma-se a categoria $C(\pi)$ dos
conjuntos finitos $E$ sobre os quais $\pi$ atua continuamente, e o funtor
identidade $F(\pi)(E)=E$ desta categoria na categoria dos conjuntos
finitos. Além disso, o grupo $\pi$ está determinado salvo isomorfismo canônico pela
condição de que $(C(\pi),F(\pi))$ seja isomorfo a um par dado. Nesse caso,
$\pi$ chama-se \emph{grupo fundamental} do esquema conexo $X$ no ponto
geométrico $a$, e denota-se por $\pi_1(X,a)$. Se $X$ não é conexo, substitui-se
o esquema pela componente conexa de $x=|a|$. Se, pelo contrário, $X$ é
conexo, os grupos $\pi_1(X,a)$ correspondentes a dois pontos geométricos
$a',a''$ de $X$ são isomorfos (o isomorfismo está determinado salvo
automorfismos interiores), de modo que, como de costume, pode-se escolher $a$
da maneira mais conveniente, por exemplo no ponto genérico de $X$, suposto
irredutível. Naturalmente, $\pi_1(X,a)$ é um funtor covariante do esquema
pontuado $(X,a)$. Todo enunciado relativo à classificação dos
recobrimentos não ramificados\footnote{O texto impresso diz «recobrimentos
inseparáveis»; o parágrafo trata da categoria $R(X)$ dos recobrimentos não
ramificados e de sua equivalência com os conjuntos finitos dotados de uma ação
contínua de $\pi_1(X,a)$.} pode ser traduzido então para a linguagem da teoria de
grupos seguindo o dicionário bem conhecido (salvo que se deve levar em conta
que aqui se trata de grupos topológicos).

\src{211}
Nosso objetivo é desenvolver o análogo da sequência exata de homotopia
dos espaços fibrados para um morfismo próprio $f:X\longrightarrow Y$.
Evidentemente, como não sabemos quais são os grupos de homotopia superiores, só
obteremos resultados necessariamente incompletos. Para poder aplicar os
teoremas fundamentais do n\textsuperscript{o} 3, devemos explicitar primeiro
alguns lemas elementares relativos aos esquemas sobre corpos ou anéis de
Artin (conforme o procedimento geral!).

\result{LEMA}{1} Seja $(X',a')$ um recobrimento não ramificado pontuado associado
a uma representação pontuada de $\pi_1(X,a)$ em um conjunto finito $E$
(provido do ponto marcado $e$). Então o morfismo canônico
$\pi_1(X',a')\longrightarrow\pi_1(X,a)$ identifica o primeiro grupo com o
estabilizador de $e$ em $\pi_1(X,a)$ (e é, portanto, injetivo).

\result{LEMA}{2} Seja $X$ um esquema algébrico sobre o corpo $k$, e seja $k'$
uma extensão radicial de $k$. Então todo recobrimento não ramificado de
$X\otimes_k k'$ provém por imagem inversa (isto é, por extensão de
escalares) de um recobrimento não ramificado de $X$, determinado salvo
isomorfismo canônico.

Desses dois lemas resulta, em particular, que para toda extensão algébrica
$K$ de $k$, e para todo ponto geométrico $a'$ de $X'=X\otimes_k K$ que se
projeta sobre o ponto geométrico $a$ de $X$, o homomorfismo funtorial
$\pi_1(X',a')\longrightarrow\pi_1(X,a)$ é injetivo.

\result{LEMA}{3} Seja $X$ um esquema completo sobre um anel local artiniano
$A$, tal que $H^0(X,\Osh_X)=A$; seja $X'$ um recobrimento não ramificado de $X$,
e seja $A'=H^0(X',\Osh_{X'})$, que é, portanto, um anel finito sobre $A$
(que \textit{a priori} pode estar ramificado sobre $A$). Sejam $X_0,X'_0$ os
subesquemas reduzidos associados a $X,X'$ (obtidos dividindo pelos feixes de
elementos nilpotentes de $\Osh_X,\Osh_{X'}$), e seja $k$ um subcorpo de
$A/\mathfrak r(A)$ sobre o qual $A/\mathfrak r(A)$ seja finito (assim, $X_0$ é
um esquema algébrico completo sobre $k$, e $X'_0$ é um recobrimento não
ramificado seu). Seja, por fim, $\Omega$ uma extensão algebricamente
fechada de $k$, e consideremos o recobrimento não ramificado
$X'_0\otimes_k\Omega$ de $X_0\otimes_k\Omega$. Então as duas condições
seguintes são equivalentes:
\begin{enumerate}
  \item[\textit{i.}] $X'_0\otimes_k\Omega$ está completamente decomposto
    sobre $X_0\otimes_k\Omega$.
  \item[\textit{ii.}] O morfismo natural
    $X'\longrightarrow X\otimes_A A'$ é um isomorfismo.
\end{enumerate}
Sob essas condições, $A'$ é uma extensão não ramificada de $A$. Por fim,
se $X'$ é conexo, a condição \textit{(i)} equivale à condição seguinte,
de aparência mais fraca:
\begin{enumerate}
  \item[\textit{i bis.}] $X'_0\otimes_k\Omega$ admite uma seção regular
    sobre $X_0\otimes_k\Omega$.
\end{enumerate}
Quando se satisfaz a condição \textit{(ii)}, diremos que o recobrimento não
ramificado $X'$ de $X$ é \emph{geometricamente trivial}.

\src{212}
\result{LEMA}{4} Seja $f:X\longrightarrow Y$ um morfismo próprio tal que
$f_*(\Osh_X)=\Osh_Y$. Sejam $a$ um ponto geométrico de $X$ e $b$ sua projeção
sobre $Y$. Então $\pi_1(X,a)\longrightarrow\pi_1(Y,b)$ é sobrejetivo.

Com efeito, é preciso demonstrar o seguinte: se um recobrimento não ramificado
$Y'$ de $Y$ (correspondente a um feixe de álgebras localmente livre
$\mathcal A$) é tal que $X\otimes_Y Y'$ é desconexo, então $Y'$ também o
é. Com efeito, $f^*(\mathcal A)$ será soma direta de dois feixes de anéis não
nulos e, portanto, também o será sua imagem direta; ora, esta última
não é outra que $\mathcal A\otimes f_*(\Osh_X)=\mathcal A$.

\result{LEMA}{5} Seja $X$ um esquema completo sobre um corpo $k$; supõe-se que
$H^0(X,\Osh_X)$ é um anel local $A$ e que $A/\mathfrak r(A)$ é radicial
sobre $k$. Seja $\Omega$ um fecho algébrico de $k$, e seja
$\overline X=X\otimes_k\Omega$ (é conexo). Escolhamos um ponto geométrico
$\overline a$ de $\overline X$ que se projete sobre o ponto geométrico $a$ de
$X$; então tem-se uma sequência exata
\[
  e\longrightarrow\pi_1(\overline X,\overline a)
  \longrightarrow\pi_1(X,a)
  \longrightarrow\pi_1(k,b)\longrightarrow e
\]
(onde $\pi_1(k,b)$ é o grupo de Galois de $\Omega$ sobre $k$).

A injetividade do primeiro homomorfismo foi vista com os lemas 1 e 2, a
exatidão no meio resulta do lema 3 e, por fim, a sobrejetividade do
último homomorfismo (que é a única que utiliza o fato de que
$A/\mathfrak r(A)$ seja radicial) resulta do lema 4.

\result{PROPOSIÇÃO}{4} Seja $f:X\longrightarrow Y$ um morfismo próprio e plano
tal que, para todo $y\in Y$, $H^0(f^{-1}(y),\Osh_{f^{-1}(y)})$ seja uma álgebra
separável sobre o corpo residual $\kappa(y)$ (o que ocorre, por exemplo, se
$f^{-1}(y)$ é um esquema separável sobre $\kappa(y)$, isto é, reduzido e tal
que os corpos correspondentes a suas componentes irredutíveis sejam
extensões separáveis de $\kappa(y)$). Então o recobrimento $Y'$ de $Y$
associado a $f_*(\Osh_X)$ é não ramificado.

A demonstração é fácil graças ao teorema 2.

\src{213}
Essa proposição, junto com o lema 1, reduz praticamente o estudo
homotópico dos morfismos próprios e planos (com fibras separáveis) ao caso em
que $f_*(\Osh_X)=\Osh_Y$ (pois, utilizando a fatoração de Stein, se
substitui $Y$ por $Y'$).

\result{OBSERVAÇÃO}{} Um morfismo plano de tipo finito cujas fibras são
esquemas separáveis (resp. lisos) chama-se \emph{separável} (resp.
\emph{liso}). Demonstra-se que, se $f$ é plano e $f^{-1}(y)$ é separável
(resp. liso), existe uma vizinhança de $f^{-1}(y)$ na qual $f$ é separável
(resp. liso). O mesmo resultado é válido para «absolutamente normal»
(«teorema de Bertini»).

Seja $f:X\longrightarrow Y$ um morfismo próprio tal que
\begin{equation}
  \tag{i}
  f_*(\Osh_X)=\Osh_Y,
\end{equation}
e seja $X'$ um esquema finito sobre $X$. Seja $Y'$ o recobrimento de $Y$
correspondente à fatoração de Stein de $X'\longrightarrow Y$
(cf. teorema 5). Seja $y\in Y$; então o conjunto das componentes conexas
da fibra $F'$ de $X'$ sobre $y$ identifica-se com o conjunto dos pontos
$y'$ de $Y'$ sobre $y$ (teorema 5). Consideremos o morfismo evidente
\begin{equation}
  \tag{*}
  X'\longrightarrow X\times_Y Y'
\end{equation}
deduzido dos morfismos naturais $X'\longrightarrow X$ e
$X'\longrightarrow Y'$\footnote{O texto impresso omite a prima no segundo
$X'$; o morfismo para $Y'$ é o da fatoração de Stein de
$X'\to Y$.}; será um isomorfismo sempre que $X'$ seja da forma
$X\times_Y Y''$, onde $Y''$ é um recobrimento não ramificado de $Y$, e
então $Y'$ não será outro senão $Y''$ e (*) será a identidade. Queremos dar
precisamente condições sob as quais $X'$ seja da forma que acabamos de
indicar, isto é, que $Y'$ seja não ramificado e (*) um isomorfismo. Para isso,
introduzamos a fibra $F$ de $X$ em $y$; é um esquema próprio sobre
$\kappa(y)$ do qual $F'$ é um recobrimento (não ramificado se $X'$ o for).
Seja $F'_1$ uma componente conexa de $F'$ correspondente a um ponto $y'_1$ de
$Y'$ sobre $y$. Suponhamos além disso
\begin{enumerate}
  \item[\textit{(ii)}] $X'$ é não ramificado sobre $X$ nos pontos de $F'_1$
    (portanto, $F'_1$ é um recobrimento não ramificado de $F$), e
  \item[\textit{(iii)}] $F'_1$ é um recobrimento geometricamente trivial de
    $F$ (cf. lema 3).
\end{enumerate}

\result{TEOREMA}{11} Nessas condições, existe uma vizinhança aberta $U'$ de
$y'_1$ em $Y'$ tal que (*) é um isomorfismo sobre $U'$. Além disso, $Y'$ é não
ramificado em $y'_1$ sobre $Y$ (mas pode ser ramificado nos demais pontos
$y'$ de $Y'$ sobre $y$).

Naturalmente, as condições \textit{(ii)} e \textit{(iii)} são também
necessárias para que seja válida a conclusão do teorema. A demonstração do
teorema é fácil, com auxílio do lema 3 e do teorema 2.

\src{214}
\result{COROLÁRIO}{1} Suponhamos sempre satisfeita \textit{(i)}. Para que um
recobrimento não ramificado sobre $X$ seja isomorfo à imagem inversa de um
recobrimento não ramificado $Y'$ de $Y$, é necessário e suficiente que $X'$
induza sobre cada fibra $f^{-1}(y)$ um recobrimento geometricamente trivial.

Como, segundo o teorema 11, o conjunto dos pontos de $Y$ para os quais se
cumpre essa condição é aberto, basta verificá-la nos pontos $y$ que são
fechados\ldots{} Assinalemos o enunciado seguinte, equivalente ao corolário 1:
o núcleo do homomorfismo (sobrejetivo segundo o lema 4)
$\pi_1(X)\longrightarrow\pi_1(Y)$ é o subgrupo invariante fechado gerado
pelas imagens em $\pi_1(X)$ dos $\pi_1(\overline{f^{-1}(y)})$, onde
\[
  \overline{f^{-1}(y)}
  =f^{-1}(y)\otimes_{\kappa(y)}\overline{\kappa(y)},
\]
$\overline{\kappa(y)}$ designa um fecho algébrico de $\kappa(y)$. Observar-se-á
que, como não se pode escolher o mesmo ponto base para todas as fibras,
os homomorfismos $\pi_1(\overline{f^{-1}(y)})\longrightarrow\pi_1(X)$ só
ficam determinados em qualquer caso (uma vez escolhido um ponto base para $X$
e, portanto, para $Y$) módulo composição por um automorfismo interior em
$\pi_1(X)$.

\result{COROLÁRIO}{2} Nas condições gerais do teorema 11, suponhamos
além disso que $Y,X,X'$ são íntegros, e sejam $K,L,L'$ seus corpos. Então existe
uma subextensão separável $K'$ de $K$ em $L'$, linearmente disjunta de $L$, tal
que $L'=LK'$ (de onde $L'=L\otimes_K K'$).

(Aplica-se a última parte do lema 3 à fibra geométrica de $X$). O caso de
aplicação mais interessante do teorema 11 obtém-se quando $f$ é um morfismo
separável. Então $X'$ é também separável sobre $Y$; portanto, em virtude
da proposição 4, $Y'$ é não ramificado sobre $Y$, e o segundo membro
$X\times_Y Y'$ de (*) é não ramificado sobre $X$. Conclui-se facilmente:

\result{COROLÁRIO}{3} Suponhamos, além de \textit{(i)}, que $f$ é separável.
Seja $X'$ um recobrimento não ramificado e conexo de $X$; para que
$X'$\footnote{O texto impresso diz aqui «$X$»; o objeto do critério é o
recobrimento $X'$ introduzido na frase anterior.} seja imagem inversa de um
recobrimento não ramificado $Y'$ de $Y$, é necessário e suficiente que o
recobrimento $\overline F'$ induzido sobre uma fibra geométrica
$\overline F=f^{-1}(y)$ admita uma seção regular.

Observar-se-á que não era necessário supor que $\overline F'$ fosse
geometricamente trivial sobre $\overline F$ (o que será verdadeiro \textit{a
posteriori}, embora \textit{a priori} esta condição seja muito mais forte). O
corolário 3 é, de fato, equivalente ao enunciado seguinte:

\src{215}
\result{COROLÁRIO}{4} Seja $f:X\longrightarrow Y$ um morfismo próprio e separável
tal que $f_*(\Osh_X)=\Osh_Y$, seja $\overline F$ a fibra geométrica de um ponto
$y\in Y$, e escolhamos um ponto geométrico em $\overline F$, do qual se obtêm,
mediante os morfismos $\overline F\longrightarrow X\longrightarrow Y$, pontos
geométricos em $X,Y$, que serão tomados como pontos base para os grupos
fundamentais de $\overline F,X,Y$. Nessas condições tem-se a sequência
exata
\[
  \boxed{
    \pi_1(\overline F)\longrightarrow\pi_1(X)
    \longrightarrow\pi_1(Y)\longrightarrow e
  }.
\]

Daí se deduzem facilmente os dois enunciados seguintes de Serre--Lang,
liberados aqui de toda hipótese de normalidade:

\result{COROLÁRIO}{5} Sejam $X,Y$ esquemas conexos sobre um corpo $k$;
suponhamos que o esquema reduzido $X_{\mathrm{réd}}$ é separável sobre $k$
(o que é automaticamente verdadeiro se $k$ é perfeito), e que $X$ ou $Y$ é
próprio sobre $k$.\footnote{Condição acrescentada na errata de 1962.}
Escolhamos um ponto geométrico $a$ (resp. $b$) em $X$ (resp. $Y$), de onde se
obtêm um ponto geométrico $c=(a,b)$ em $X\times_kY$ e um morfismo natural
\[
  \pi_1(X\times_kY,c)\longrightarrow
  \pi_1(X,a)\times\pi_1(Y,b)
\]
(deduzido dos morfismos funtoriais de $\pi_1(X\times Y,c)$ em $\pi_1(X,a)$
e $\pi_1(Y,b)$). Este morfismo é injetivo, e inclusive bijetivo se $k$ é
algebricamente fechado.

(A sobrejetividade neste último caso é quase trivial). Daí se deduz,
com Serre--Lang:

\result{COROLÁRIO}{6} Seja $X$ um esquema algébrico conexo sobre um corpo
algebricamente fechado $k$, e seja $K$ uma extensão algebricamente fechada
de $k$; então os grupos fundamentais de $X$ e $X\otimes_kK$ são os
mesmos, isto é, todo recobrimento não ramificado deste último esquema
provém por extensão de escalares de um recobrimento não ramificado (único
salvo isomorfismo) de $X$.

\result{OBSERVAÇÕES}{}

1) Utilizando a proposição 4, vê-se que a hipótese
$f_*(\Osh_X)=\Osh_Y$ do corolário 4 não é essencial. No caso geral, em vez
de colocar o grupo unidade $e$ depois de $\pi_1(Y)$, é preciso continuar com
\[
  \pi_0(\overline F)\longrightarrow\pi_0(X)
  \longrightarrow\pi_0(Y)\longrightarrow e.
\]
como em topologia algébrica.

\src{216}
2) Em geral, por ora não se pode dizer nada sobre o núcleo de
$\pi_1(\overline F)\longrightarrow\pi_1(X)$, no qual deveria intervir um
$\pi_2(Y)$. Parece, no entanto, que se deveria poder demonstrar que
$\pi_1(\overline F)\longrightarrow\pi_1(X)$ é injetivo se $Y$ é o espectro
de um anel local $A$, apoiando-se no teorema 12 mais adiante (que afirma que
isso ocorre se $A$ é completo).

O teorema 11 utilizava somente os teoremas 1 e 2. Vamos aplicar agora o
teorema 3, apoiando-nos no lema elementar seguinte:

\result{LEMA}{6} Seja $X$ um esquema, e $X_0$ o esquema reduzido correspondente
(isto é, no qual se anularam os elementos nilpotentes). Então todo
recobrimento não ramificado $X'_0$ de $X_0$ é induzido por um recobrimento
não ramificado $X'$ de $X$, determinado salvo isomorfismo canônico.

Este lema, de natureza puramente local, desempenha aqui um papel análogo ao do
teorema 8 em teoria de módulos. Combinando-o com o teorema de existência
(teorema 3), deduz-se:

\result{TEOREMA}{12} Seja $A$ um anel local noetheriano e completo de corpo
residual $k$. Seja $X$ um esquema próprio sobre $A$. Então todo recobrimento
não ramificado $X'_0$ de $X_0=X\otimes_A k$ é induzido por um recobrimento não
ramificado $X'$ de $X$, único salvo isomorfismo.

Em outros termos:

\result{COROLÁRIO}{1} Escolhamos um ponto geométrico em $X_0$ como ponto base para
os grupos fundamentais de $X_0$ e $X$. Então o homomorfismo canônico
$\pi_1(X_0)\longrightarrow\pi_1(X)$ é um isomorfismo.

Apliquemos agora a $X_0$ o lema 5 (supondo, para simplificar, que
$H^0(X_0,\Osh_{X_0})=k$), e observemos que, por ser $A$ completo, suas
extensões não ramificadas correspondem às extensões não ramificadas de seu
corpo residual, isto é, tem-se $\pi_1(Y)=\pi_1(k)$ (onde
$Y=\operatorname{Spec}(A)$). Obtém-se a sequência exata:
\[
  e\longrightarrow\pi_1(\overline X_0)
  \longrightarrow\pi_1(X)\longrightarrow\pi_1(Y)\longrightarrow e.
\]

\result{COROLÁRIO}{2} Seja $f:X\longrightarrow Y$ um morfismo próprio e plano,
sejam $y_1$ um ponto de $Y$ e $y_0$ uma especialização de $y_1$, e consideremos
as fibras «geométricas» correspondentes $\overline X_1$ e $\overline X_0$;
supõe-se que $\overline X_0$ é separável e conexa (o que implica que
$\overline X_1$ satisfaz as mesmas condições). Então pode-se encontrar um
homomorfismo de grupos
$\pi_1(\overline X_1)\longrightarrow\pi_1(\overline X_0)$, definido salvo um
automorfismo interior, e este homomorfismo é sobrejetivo.

\src{217}
Poder-se-ia esperar que este homomorfismo fosse sempre bijetivo; infelizmente,
em geral não é assim se $\kappa(y_0)$ tem característica $>0$. Obteremos,
no entanto, mais adiante uma cota para o núcleo deste homomorfismo (ao menos se
$\overline X_0$ é liso), da qual se segue que, se $\kappa(y_0)$ tem
característica $0$, o homomorfismo anterior é bijetivo (resultado que também
poderia ser demonstrado por via transcendente). Em todo caso, obtemos já uma cota
para o $\pi_1$ de uma fibra especial por meio daquele de uma fibra genérica.
Utilizando, por exemplo, o fato de que uma curva algébrica em característica
$p$ se pode levantar a uma curva em característica $0$ (teorema 9, corolário
4), obtém-se por via transcendente:

\result{COROLÁRIO}{3} Seja $X_0$ o esquema de uma curva completa lisa sobre
um corpo algebricamente fechado de característica arbitrária, e seja $g$ o
gênero de $X_0$; então $\pi_1(X_0)$ admite $2g$ geradores topológicos,
ligados pela relação bem conhecida.

Daí se deduz, mediante uma técnica bem conhecida de seções
hiperplanas:

\result{COROLÁRIO}{4} Seja $X$ um esquema projetivo liso sobre um corpo
algebricamente fechado de característica arbitrária; então $\pi_1(X)$
admite um número finito de geradores topológicos.

Determinemos o núcleo do homomorfismo
$\pi_1(\overline X_1)\longrightarrow\pi_1(\overline X_0)$. Para isso, pode-se
supor que $Y$ é o espectro de um anel de valoração discreta completo
$V=\Osh_y$ ($y=y_0$). A questão equivale à seguinte: dado um recobrimento
não ramificado $\overline X'_1$ de $\overline X_1$ (que, se se quiser, pode
ser suposto galoisiano), determinar em que condições provém de um
recobrimento não ramificado de $X_0$. \textit{A priori}, o recobrimento dado
provém, por extensão de escalares, de um recobrimento não ramificado $X'_1$
de $X_1\otimes_K K'$, onde $K'$ é uma subextensão finita do fecho
algébrico $\overline K$ do corpo de frações $K$ de $V$; se
$\overline X'_1$ fosse galoisiano de grupo $G$, pode-se escolher $X'_1$
galoisiano de grupo $G$. Tem-se então: para que
$X'_1\otimes_{K'}\overline K=\overline X'_1$ provenha de um recobrimento não
ramificado de $X_0$, é necessário e suficiente que se possa encontrar uma extensão
finita $K''$ de $K'$ em $\overline K$ tal que
$X''_1=X'_1\otimes_{K'}K''$ seja da forma $X''\otimes_{V''}K''$, onde $V''$
é o fecho normal de $V$ em $K''$ e $X''$ é um recobrimento não ramificado
de $X\otimes_VV''$. Suponhamos, por exemplo, que $X_0$ é absolutamente normal;
daí resulta que $X\otimes_VV''$ é normal (pois é plano sobre $V''$ e tem
fibra especial normal), e seu corpo de funções é idêntico ao $K''(X_1)$ de
\[
  X_1\otimes_KK''
  =(X\otimes_VK)\otimes_KK''
  =X\otimes_VK''
  =(X\otimes_VV'')\otimes_{V''}K''.
\]
Seja $L''=K''(X'_1)$ o corpo de funções de $X'_1\otimes_{K'}K''$; é uma
extensão finita separável de $K''(X_1)$, e a condição anterior significa
também que $L''$ é uma extensão não ramificada do corpo de funções de
$X\otimes_VV''$ (isto é, a normalização de $X\otimes_VV''$ em $L''$ é não
ramificada sobre $X\otimes_VV''$). Basta, de resto, verificar que $L''$ é
não ramificado nos pontos da fibra especial de $X\otimes_VV''$ (pois é não
ramificado sobre a fibra genérica $X_1\otimes_KK''$). Se agora $X_0$ é liso,
do «teorema de pureza» de Nagata--Zariski resulta inclusive que basta verificar
que $L''$ é não ramificado sobre o anel local $O''$ do ponto genérico da
fibra especial de $X\otimes_VV''$, que é um anel de valoração discreta,
normalização em $K''(X_1)$ do anel local $O\subset K(X_1)$ do ponto
genérico da fibra especial de $X$. Portanto, obteve-se:

\src{218}
\result{COROLÁRIO}{5} Nas condições e com as notações anteriores, para
que o recobrimento não ramificado
$X'_1\otimes_{K'}\overline K$ de
$\overline X_1=X_1\otimes_K\overline K$ provenha de um recobrimento não
ramificado de $\overline X_0$, é necessário e suficiente que exista uma
subextensão finita $K''$ de $\overline K/K'$ tal que $K''(X'_1)$ seja não
ramificado sobre o anel de valoração discreta $O''\subset K''(X_1)$.

Observemos agora que $O''$ é a normalização em $K''(X_1)$ do anel de
valoração discreta $O'\subset K'(X_1)$ (normalização de $O$ em $K'(X_1)$), e
que $O'$ contém a normalização $V'$ de $V$ em $K'$; um uniformizante $u$ de
$V'$ é também um uniformizante de $O'$. Suponhamos agora que $X'_1$ é
galoisiano de grupo de Galois $G$, de ordem $n$ primo com a característica $p$
de $\kappa(y_0)$ (que é também a característica do corpo residual de
$O'$). Portanto, $K'(X'_1)$ está «moderadamente ramificado» sobre $O'$, de onde
resulta facilmente («lema de Abhyankar») que, se se lhe adjunta uma raiz
$n$-ésima $v$ de um uniformizante de $O'$, torna-se não ramificado sobre a
normalização de $O'$ em $K'(X_1)(v)$. Ora, pode-se tomar como $v$ uma
raiz $n$-ésima de um uniformizante de $V'$, o que prova que se satisfaz a
condição do corolário 5. (Essa possibilidade de utilizar o lema de Abhyankar e
o teorema de pureza foi-me vendida por Serre). Para exprimir o resultado
obtido, introduzamos, para todo grupo compacto totalmente descontínuo $\pi$,
o grupo quociente $\widetilde\pi$ de $\pi$ pelo subgrupo fechado gerado
por seus $p$-subgrupos de Sylow, isto é, o limite projetivo dos grupos
quocientes discretos de $\pi$ que são de ordem prima com $p$. Com esta notação,
obtém-se:

\result{TEOREMA}{13} Seja $f:X\longrightarrow Y$ um morfismo próprio e plano,
sejam $y_1$ um ponto de $Y$ e $y_0$ uma especialização de $y_1$, e suponhamos
$\overline X_0$ conexo e liso. Então o homomorfismo
$\widetilde\pi_1(\overline X_1)\longrightarrow
\widetilde\pi_1(\overline X_0)$ deduzido do homomorfismo sobrejetivo do
teorema 12, corolário 2, é um isomorfismo.

Em outros termos:

\src{219}
\result{COROLÁRIO}{1} A classificação dos recobrimentos galoisianos não
ramificados, de grupo de Galois de ordem prima com a característica $p$ de
$\kappa(y_0)$, é a mesma para $\overline X_0$ e para $\overline X_1$.

Em particular, se $\kappa(y_0)$ tem característica zero, obtém-se por via
algébrica que $\pi_1(\overline X_1)\longrightarrow\pi_1(\overline X_0)$ é
bijetivo.

Assinalemos, por fim, que as técnicas utilizadas proporcionam também o
resultado seguinte, mais geral que o teorema 13:

\result{TEOREMA}{14} Seja $f:X\longrightarrow Y$ um morfismo próprio e liso,
seja $D$ um subesquema fechado de $X$, liso sobre $Y$, de codimensão 1 em
todos os seus pontos. Para uma fibra $Z=f^{-1}(z)$ de $f$, seja $Z'=Z-Z\cap D$ e seja
$\pi_1^t(Z')$ o quociente do grupo fundamental $\pi_1(Z')$ que classifica os
recobrimentos não ramificados de $Z'$ que estão «moderadamente ramificados»
sobre $\overline{Z\cap D}$. Sejam $y_0,y_1$ como no teorema 13. Então se
tem um homomorfismo sobrejetivo (definido salvo um automorfismo interior)
\[
  \pi_1^t(\overline X'_1)\longrightarrow\pi_1^t(\overline X'_0),
\]
e o homomorfismo correspondente
\[
  \widetilde\pi_1^t(\overline X'_1)
  \longrightarrow\widetilde\pi_1^t(\overline X'_0)
\]
é um isomorfismo.

Daí se deduzem as variantes correspondentes dos corolários do
teorema 13 e do corolário 4 do teorema 12. Do mesmo modo, utilizando o
teorema 9, corolário 3, obtém-se por via transcendente:

\result{COROLÁRIO}{} Seja $X_0$ o esquema de uma curva completa lisa sobre um
corpo algebricamente fechado de característica arbitrária, e seja
$S=(s_i)_{1\le i\le n}$ uma parte finita com $n$ elementos de $X_0$. Então
$\pi_1^t(X_0-S)$ admite $2g+n$ geradores topológicos
$x_i,y_i$ ($1\le i\le g$) e $\sigma_j$ ($1\le j\le n$), ligados pela relação
\[
  \left(\prod_i x_i y_i x_i^{-1}y_i^{-1}\right)
  \sigma_1\cdots\sigma_n=1,
\]
onde os $\sigma_j$ são geradores dos grupos de inércia correspondentes
aos $s_j$. Para todo grupo finito $G$ de ordem prima com a característica,
gerado por elementos $x_i,y_i,\sigma_j$ que satisfazem a relação anterior,
existe um recobrimento galoisiano não ramificado de $X_0-S$, de grupo $G$, cujos
grupos de inércia nos pontos $s_j$ estão gerados pelos $\sigma_j$.

Quando $X_0$ é de gênero 0 e $n=3$, tem-se uma solução do «problema dos
três pontos», ao menos para os recobrimentos galoisianos de ordem prima com
a característica. (Aqui, de resto, o teorema 9 é inútil, e parece, por
outra parte, que o corolário anterior pode ser deduzido do caso particular
considerado no problema dos três pontos).

\src{220}
\result{OBSERVAÇÕES}{}

1) Um estudo mais completo, no qual sem dúvida interviriam recobrimentos
galoisianos generalizados de $X,X_0,X_1$ (com grupo de Galois eventualmente
infinitesimal), deveria permitir recuperar o núcleo do teorema 12, corolário
2. Em contrapartida, o estudo dos recobrimentos que admitem ramificação não
«moderada» parece ter de ser muito mais difícil.

2) O lema 6, junto com um resultado de Grauert relativo ao completamento
formal de um esquema projetivo não singular ao longo de uma seção
hiperplana (ou com o teorema, ainda não demonstrado, mencionado na
observação 2 depois do teorema 11), permite também demonstrar em geometria
algébrica «abstrata» o teorema clássico de Lefschetz sobre o grupo
fundamental.

\begin{center}
  \textbf{BIBLIOGRAFIA}
\end{center}

\begin{description}
  \item[{[1]}] Dieudonné (J.) e Grothendieck (A.). --- \emph{Éléments de
    géométrie algébrique}, Publications mathématiques de l'Institut des Hautes
    Études Scientifiques (de próxima publicação).
  \item[{[2]}] Grothendieck (Alexander). --- \emph{The cohomology theory of
    abstract algebraic varieties}. International Congress of Mathematicians
    [1958. Edinburgh] (de próxima publicação).
  \item[{[3]}] Serre (Jean-Pierre). --- \emph{Faisceaux algébriques
    cohérents}, Annals of Math., Series 2, t. 61, 1955, p. 197--278.
  \item[{[4]}] Serre (Jean-Pierre). --- \emph{Géométrie algébrique et
    géométrie analytique}, Ann. Institut Fourier Grenoble, t. 6, 1955--1956,
    p. 1--42.
\end{description}
